% Solutions to the exercises of Chapter 20.
\begin{solutions}
  \item From $1/z_i = 1/f_L - 1/z_o$, $z_i = f_Lz_o/(z_o - f_L) = 4\cdot500/496$~mm $= 4.0323$~mm: the sensor must move \SI{32}{\micro\metre} away from the lens.
        To first order the shift is $f_L^2/z_o$, so it shrinks with the square of the focal length. If the sensor is left at $f_L$, the light from the point
        converges \SI{32}{\micro\metre} behind it, and an aperture of diameter $D$ produces a blur spot of diameter $D(z_i - f_L)/z_i$ on the sensor. For an f/2
        lens ($D = 2$~mm) this is \SI{16}{\micro\metre}, under three pixels of \SI{6}{\micro\metre}. Short focal lengths give a large depth of field, which is why small
        fixed-focus cameras stay nearly sharp over a wide range of distances.

  \item The principal point is not at the image center, so add the angles from the optical axis to the two opposite edges, measured along the row and
        the column through $(c_u, c_v)$. Without distortion the left and right edges are at normalized $x = -367.7/458.65$ and $384.3/458.65$, at
        $38.7^\circ$ and $40.0^\circ$ from the axis; the top and bottom edges are at $28.6^\circ$ and $26.8^\circ$. With distortion, the edge pixels must first be
        undistorted (Exercise~5): the middle of the left edge, distorted coordinate $x_d = -0.80$, is the image of a ray with $x = -1.02$, at $45.6^\circ$. Altogether:
        \begin{center}\small
        \begin{tabular}{@{}l c c c@{}}
        \toprule
        & horizontal & vertical & diagonal \\
        \midrule
        pinhole only, $2\arctan(W/2f_u)$ & $78.7^\circ$ & $55.4^\circ$ & $88.4^\circ$ \\
        with distortion & $93.1^\circ$ & $59.7^\circ$ & $106.2^\circ$ \\
        \bottomrule
        \end{tabular}
        \end{center}
        Barrel distortion compresses the periphery onto the sensor, so the camera sees $14^\circ$ more horizontally than its focal length suggests; the image
        corners look along rays $53^\circ$ off the axis, not $44^\circ$.

  \item The lines are $l_1 = (1, 1, -1)$ and $l_2 = (1, -1, 0)$, and $l_1\times l_2 = (-1, -1, -2) \sim (\tfrac12, \tfrac12, 1)$: the intersection is $(\tfrac12, \tfrac12)$, which
        satisfies both equations. The line through the points is $(1, 2, 1)\times(3, 1, 1) = (1, 2, -5)$, i.e.\ $x + 2y = 5$; both points satisfy it.

  \item For parallel lines $(a, b, c_1)$ and $(a, b, c_2)$, $c_1 \ne c_2$,
        \[
          (a, b, c_1)\times(a, b, c_2) = \big(b(c_2 - c_1),\ a(c_1 - c_2),\ 0\big) = (c_2 - c_1)\,(b, -a, 0),
        \]
        a point with third coordinate zero, i.e.\ at infinity, in the direction $(b, -a)$ of the lines themselves (their normal is $(a, b)$). For the 3D lines
        $X(\lambda) = X_0 + \lambda d$, the camera $P = K[I\ \ 0]$ gives the image $K(X_0 + \lambda d) \sim K(X_0/\lambda + d)$, which tends to
        \[
          \tilde v = Kd \qquad (\lambda\to\infty),
        \]
        independently of $X_0$: all lines of the family pass through this vanishing point, at pixel $(f_ud_x/d_z + c_u,\ f_vd_y/d_z + c_v)$ for zero skew. If $d_z = 0$ (lines
        parallel to the image plane), $Kd$ has third coordinate zero and the image lines stay parallel. The vanishing points of all directions in a plane with
        normal $n$ lie on the line $K\mT n$, the plane's horizon, since $(K\mT n)\T Kd = n\T d = 0$.

  \item Distortion and the fixed-point iteration are the formulas of \cref{sec:distortion}. For Gauss--Newton,
        solve $f(x) = x_d$ with the $2\times2$ Jacobian $J_f$ of the distortion: $x \gets x - J_f\inv\big(f(x) - x_d\big)$ (with as many equations as unknowns, this is
        Newton's method). Measuring the error against a fully converged solution, in normalized units ($10^{-6}$ is $5\times10^{-4}$~px):
        \begin{center}\small
        \begin{tabular}{@{}l c c c@{}}
        \toprule
        pixel & contraction factor & fixed point & Gauss--Newton \\
        \midrule
        the four image corners & $0.03$--$0.17$ & 5--8 iterations & 2--3 iterations \\
        middle of left edge, $(0, c_v)$ & $0.34$ & 12 iterations & 3 iterations \\
        \bottomrule
        \end{tabular}
        \end{center}
        The fixed-point error falls by a constant factor per iteration. For radial distortion this factor is $|2k_1r^2 + 4k_2r^4|/(1 + k_1r^2 + k_2r^4)$, which peaks at
        $0.34$ for $r \approx 1$: the hardest point is the middle of the left and right edges, not the corners, where the $k_2$ term flattens the
        distortion curve. Gauss--Newton converges quadratically (errors $10^{-2}, 10^{-4}, 10^{-8}$ at a corner), and each iteration costs only a few more
        operations. A practical warning: OpenCV's \code{cv::undistortPoints} stops after five iterations unless it is given termination criteria, which
        leaves an error of about $1.1\times10^{-3}$, half a pixel, at the middle of the left edge. Milestone 8's round-trip test catches this.

  \item The homography is $H = [h_1\ h_2\ h_3] = \lambda K[r_1\ r_2\ t]$ for an unknown scale $\lambda$, so $r_1 = \lambda\inv K\inv h_1$ and $r_2 = \lambda\inv K\inv h_2$. The columns of a
        rotation are orthonormal:
        \[
          r_1\T r_2 = 0 \;\Rightarrow\; h_1\T K\mT K\inv h_2 = h_1\T\omega h_2 = 0, \qquad
          r_1\T r_1 = r_2\T r_2 \;\Rightarrow\; h_1\T\omega h_1 = h_2\T\omega h_2 .
        \]
        The unknown $\lambda$ cancels in both; the unit-norm condition itself cannot be used, because it would involve $\lambda$. Each equation is linear in the six
        distinct entries $b = (\omega_{11}, \omega_{12}, \omega_{22}, \omega_{13}, \omega_{23}, \omega_{33})$: $h_i\T\omega h_j = v_{ij}\T b$ with
        \[
          v_{ij} = \big(h_{i1}h_{j1},\ h_{i1}h_{j2} + h_{i2}h_{j1},\ h_{i2}h_{j2},\ h_{i3}h_{j1} + h_{i1}h_{j3},\ h_{i3}h_{j2} + h_{i2}h_{j3},\ h_{i3}h_{j3}\big).
        \]
        Each image contributes the rows
        $v_{12}\T$ and $(v_{11} - v_{22})\T$, and $b$ is the null vector of the stacked matrix (the smallest right singular vector), determined up to scale. Since
        $\omega$ has five degrees of freedom up to scale, three images in general position suffice. Then $\omega = K\mT K\inv = LL\T$ with $L = K\mT$ lower
        triangular, so the Cholesky factor of $\omega$ (sign and scale fixed so that $\omega \succ 0$) gives $K = L\mT$, normalized to $K_{33} = 1$.

  \item Translation at speed $v$ perpendicular to the ray moves the image at $v/Z = 0.1$~rad/s, about 46~px/s. Rotation at $\omega$ about the camera's $y$-axis moves a point at image position $x$ at $\omega(1 + x^2) \ge \omega =
        1$~rad/s, 460~px/s, independent of depth. Rotation dominates by a factor of ten. The translational signal of a point scales as $1/Z$, so distant points
        constrain translation weakly and points at infinity not at all, while they constrain rotation as well as near points do. Worse, a rotation error of
        $\delta\omega$ produces the same image motion as a velocity error of $Z\,\delta\omega$: a gyroscope-rate error of $0.01$~rad/s mimics $0.1$~m/s at 10~m. Estimating
        translation from distant points therefore requires the rotation to be known very accurately, which is exactly what the gyroscope provides in VIO.
\end{solutions}
